\item Langkah kedua adalah (silahkan jelaskan langkah-langkah
      selama praktikum)

      Jika ada tabel yang perlu dilampirkan berikut format table nya.

      \begin{table}[H]
        \centering
        \caption{Contoh penggunaan table}
        \linespread{1}\fontsize{7}{8}\selectfont
        \rowcolors{2}{tabband}{white}
        \setlength{\tabcolsep}{3pt}
        \renewcommand{\arraystretch}{1.25}
        \arrayrulecolor{tabline}
        \begin{tabular}{!{\color{tabline}\vrule}>{\centering\arraybackslash}p{9mm}!{\color{tabline}\vrule}p{34mm}!{\color{tabline}\vrule}p{34mm}!{\color{tabline}\vrule}p{34mm}!{\color{tabline}\vrule}}
          \hline
          \rowcolor{tabhead}
          \thead{No} & \thead{Judul kolom 1} & \thead{Judul Kolom 2} & \thead{Judul Kolom 3} \\
          \hline
          & & & \\ \hline
          & & & \\ \hline
          & & & \\ \hline
          & & & \\ \hline
          & & & \\ \hline
          & & & \\ \hline
          & & & \\ \hline
          & & & \\ \hline
        \end{tabular}
      \end{table}